\section{Predicted questions and rapid-fire drill}

\subsection{Thesis questions (first 5 minutes)}

\begin{summary}{Your thesis in two minutes}
rCFD records a flow once and replays it to carry a scalar. It was built for passive scalars; I
tested it on temperature, which drives its own flow (buoyant jet in a $75\times75\times15$~mm tank,
Re 99). (1) No single mixing constant is best: the outlet sits under the front, so smearing
drains heat. (2) A bounded, vertically weighted anti-diffusive correction brings the front from
15.7 to 5.3/4.5~mm (reference 5.0), no value outside the data range. (3) Under wall cooling the
replay was 4.34~K too warm because the reference used $c_p$ of air (1006.43) and the replay of
water (4182); on one fluid 0.66~K, corrected 0.06~K on 9.62~K. (4) What remains is missing
transport: a recorded flux cuts the heterogeneity error from 1.66 to 0.26~K. Take-home: a
database belongs to a flow regime.
\end{summary}

\pic{A cooking video recorded once (preparation) and replayed to cook many times (replay). It
works for soup that just gets stirred (passive scalar). For bread that rises by itself (buoyancy),
the video does not show how your dough will move, so you must correct it.}

\begin{qabox}[chair]{\asked What problem did you solve, and why ML/AI?}
\from{Thesis, Appendix A1 (story summary) · asked in past defenses}
\ans I measured what a recurrence replay loses when the scalar is not passive, and repaired most of
it without solving the flow again. It is AI for simulation: expensive solver data is reused by a
cheap model; choosing the next frame is a nearest-neighbour search; and my main result is a
generalization result: the database covers a regime, and a new boundary condition is
out-of-distribution.
\end{qabox}

\begin{qabox}[chair]{\asked What would you do differently?}
\from{Thesis, Ch.~8 and future work · asked in past defenses}
\ans Check the material properties of both simulations on the first day: the specific-heat mismatch
cost the most time. Then an implicit exchange so the recorded diffusivity can act.
\end{qabox}

\begin{qabox}[Klambauer]{\pred How would you solve your problem with deep learning?}
\from{Thesis $\times$ DL: graph neural networks · predicted}
\ans Learn the correction $\mathcal{C}$ with a graph neural network on the mesh: input the field
and the recorded map, output the next field, loss = distance to the reference plus a penalty on the
energy balance. The risk is the same as in my thesis: it only generalizes inside the regime of its
training data.
\end{qabox}

\begin{qabox}[Hochreiter]{\pred Could an LSTM replace the replay?}
\from{Thesis $\times$ LSTM \& Transformers (Hochreiter) · predicted}
\ans An LSTM (or xLSTM) could learn the time evolution of a low-dimensional state, e.g.\ the POD
coefficients of Figure 5.9, as a sequence-to-sequence model. It would be cheap and could carry
memory of the filling phase, which is exactly where my replay lags. But it needs many trajectories
to generalize, and without a conservation constraint it can drift in energy, which rCFD
guarantees by construction.
\end{qabox}

\subsection{Rapid-fire (say it in one breath)}

\begin{center}\small
\begin{tabularx}{\linewidth}{@{}>{\raggedright\arraybackslash}p{0.40\linewidth}X@{}}
\toprule
Question & Answer \\ \midrule
Most powerful learning machine? & The brain; closest method: neural networks. \\
Delta, softmax + CE? & $\hat{\mathbf{y}}-\mathbf{y}$. \\
Activation from Linz? & SELU (self-normalizing), ELU. \\
Who found vanishing gradients? & Hochreiter, 1991, Linz. \\
Constant error carousel? & Cell self-loop weight 1, linear: error flows unchanged. \\
Forget gate added by? & Gers, Schmidhuber, Cummins, 2000. \\
GRU gates? & Update (coupled input/forget) and reset. \\
BPTT memory? & Linear in $T$; RTRL: independent of $T$, $O(N^4)$ per step. \\
Truncated BPTT? & Cuts the past after $k$ steps; approximate gradient. \\
Gradient clipping fixes? & Exploding only, not vanishing. \\
Why $\sqrt{d_k}$? & Keeps dot-product variance at 1; softmax does not saturate. \\
Attention complexity? & $O(N^2)$; RNN $O(N)$. \\
Positional encoding why? & Attention is permutation-equivariant. \\
xLSTM two ideas? & Exponential gating; matrix memory (mLSTM). \\
Hopfield $\leftrightarrow$ attention? & $X\,\mathrm{softmax}(\beta X^\top\bm\xi)$ = attention. \\
UAT hidden layers? & One, finitely many units. \\
Non-differentiable loss? & 0--1 loss; use cross-entropy or hinge. \\
Regularization? & Weight decay, dropout, early stopping, augmentation, FMS. \\
Optimizers? & SGD, momentum, AdaGrad, RMSprop, Adam. \\
Output activation, regression? & Linear + squared loss. \\
VC dim.\ of lines in 2D? & 3. \\
Kernel trick? & Dot products $\to$ $k(\mathbf{x},\mathbf{x}')$. \\
MLE problem? & Overfits with little data (sunrise problem). \\
Your Re / grid / frames? & 99 / 1~mm cubes / 600 frames. \\
Your speed-up? & About 260 times. \\
\bottomrule
\end{tabularx}
\end{center}
